%!TeX root=../emmatop.tex
\chapter[Chapter \thechapter]{}
\lettrine[lines=4,lraise=0.3]{E}{mma} was not required, by any subsequent discovery, to retract her ill opinion of Mrs Elton. Her observation had been pretty correct. Such as Mrs Elton appeared to her on this second interview, such she appeared whenever they met again,—self-important, presuming, familiar, ignorant, and ill-bred. She had a little beauty and a little accomplishment, but so little judgment that she thought herself coming with superior knowledge of the world, to enliven and improve a country neighbourhood; and conceived Miss Hawkins to have held such a place in society as Mrs Elton's consequence only could surpass.

There was no reason to suppose Mr Elton thought at all differently from his wife. He seemed not merely happy with her, but proud. He had the air of congratulating himself on having brought such a woman to Highbury, as not even Miss Woodhouse could equal; and the greater part of her new acquaintance, disposed to commend, or not in the habit of judging, following the lead of Miss Bates's good-will, or taking it for granted that the bride must be as clever and as agreeable as she professed herself, were very well satisfied; so that Mrs Elton's praise passed from one mouth to another as it ought to do, unimpeded by Miss Woodhouse, who readily continued her first contribution and talked with a good grace of her being »very pleasant and very elegantly dressed.«

In one respect Mrs Elton grew even worse than she had appeared at first. Her feelings altered towards Emma.—Offended, probably, by the little encouragement which her proposals of intimacy met with, she drew back in her turn and gradually became much more cold and distant; and though the effect was agreeable, the ill-will which produced it was necessarily increasing Emma's dislike. Her manners, too—and Mr Elton's, were unpleasant towards Harriet. They were sneering and negligent. Emma hoped it must rapidly work Harriet's cure; but the sensations which could prompt such behaviour sunk them both very much.—It was not to be doubted that poor Harriet's attachment had been an offering to conjugal unreserve, and her own share in the story, under a colouring the least favourable to her and the most soothing to him, had in all likelihood been given also. She was, of course, the object of their joint dislike.—When they had nothing else to say, it must be always easy to begin abusing Miss Woodhouse; and the enmity which they dared not shew in open disrespect to her, found a broader vent in contemptuous treatment of Harriet.

Mrs Elton took a great fancy to Jane Fairfax; and from the first. Not merely when a state of warfare with one young lady might be supposed to recommend the other, but from the very first; and she was not satisfied with expressing a natural and reasonable admiration—but without solicitation, or plea, or privilege, she must be wanting to assist and befriend her.—Before Emma had forfeited her confidence, and about the third time of their meeting, she heard all Mrs Elton's knight-errantry on the subject.—

\begin{figure}[tbph]
\centering
\includegraphics[width=.7\linewidth]{33charming}
\caption{»You have heard those charming lines of the poet«}
\end{figure}

»Jane Fairfax is absolutely charming, Miss Woodhouse.—I quite rave about Jane Fairfax.—A sweet, interesting creature. So mild and ladylike—and with such talents!—I assure you I think she has very extraordinary talents. I do not scruple to say that she plays extremely well. I know enough of music to speak decidedly on that point. Oh! she is absolutely charming! You will laugh at my warmth—but, upon my word, I talk of nothing but Jane Fairfax.—And her situation is so calculated to affect one!—Miss Woodhouse, we must exert ourselves and endeavour to do something for her. We must bring her forward. Such talent as hers must not be suffered to remain unknown.—I dare say you have heard those charming lines of the poet,

\begin{verse}
Full many a flower is born to blush unseen,\\
    And waste its fragrance on the desert air.
\end{verse}

We must not allow them to be verified in sweet Jane Fairfax.«

»I cannot think there is any danger of it,« was Emma's calm answer—»and when you are better acquainted with Miss Fairfax's situation and understand what her home has been, with Colonel and Mrs Campbell, I have no idea that you will suppose her talents can be unknown.«

»Oh! but dear Miss Woodhouse, she is now in such retirement, such obscurity, so thrown away.—Whatever advantages she may have enjoyed with the Campbells are so palpably at an end! And I think she feels it. I am sure she does. She is very timid and silent. One can see that she feels the want of encouragement. I like her the better for it. I must confess it is a recommendation to me. I am a great advocate for timidity—and I am sure one does not often meet with it.—But in those who are at all inferior, it is extremely prepossessing. Oh! I assure you, Jane Fairfax is a very delightful character, and interests me more than I can express.«

»You appear to feel a great deal—but I am not aware how you or any of Miss Fairfax's acquaintance here, any of those who have known her longer than yourself, can shew her any other attention than«—

»My dear Miss Woodhouse, a vast deal may be done by those who dare to act. You and I need not be afraid. If we set the example, many will follow it as far as they can; though all have not our situations. We have carriages to fetch and convey her home, and we live in a style which could not make the addition of Jane Fairfax, at any time, the least inconvenient.—I should be extremely displeased if Wright were to send us up such a dinner, as could make me regret having asked more than Jane Fairfax to partake of it. I have no idea of that sort of thing. It is not likely that I should, considering what I have been used to. My greatest danger, perhaps, in housekeeping, may be quite the other way, in doing too much, and being too careless of expense. Maple Grove will probably be my model more than it ought to be—for we do not at all affect to equal my brother, Mr Suckling, in income.—However, my resolution is taken as to noticing Jane Fairfax.—I shall certainly have her very often at my house, shall introduce her wherever I can, shall have musical parties to draw out her talents, and shall be constantly on the watch for an eligible situation. My acquaintance is so very extensive, that I have little doubt of hearing of something to suit her shortly.—I shall introduce her, of course, very particularly to my brother and sister when they come to us. I am sure they will like her extremely; and when she gets a little acquainted with them, her fears will completely wear off, for there really is nothing in the manners of either but what is highly conciliating.—I shall have her very often indeed while they are with me, and I dare say we shall sometimes find a seat for her in the barouche-landau in some of our exploring parties.«

»Poor Jane Fairfax!«—thought Emma.—»You have not deserved this. You may have done wrong with regard to Mr Dixon, but this is a punishment beyond what you can have merited!—The kindness and protection of Mrs Elton!—»Jane Fairfax and Jane Fairfax.« Heavens! Let me not suppose that she dares go about, Emma Woodhouse-ing me!—But upon my honour, there seems no limits to the licentiousness of that woman's tongue!«

Emma had not to listen to such paradings again—to any so exclusively addressed to herself—so disgustingly decorated with a »dear Miss Woodhouse.« The change on Mrs Elton's side soon afterwards appeared, and she was left in peace—neither forced to be the very particular friend of Mrs Elton, nor, under Mrs Elton's guidance, the very active patroness of Jane Fairfax, and only sharing with others in a general way, in knowing what was felt, what was meditated, what was done.

She looked on with some amusement.—Miss Bates's gratitude for Mrs Elton's attentions to Jane was in the first style of guileless simplicity and warmth. She was quite one of her worthies—the most amiable, affable, delightful woman—just as accomplished and condescending as Mrs Elton meant to be considered. Emma's only surprize was that Jane Fairfax should accept those attentions and tolerate Mrs Elton as she seemed to do. She heard of her walking with the Eltons, sitting with the Eltons, spending a day with the Eltons! This was astonishing!—She could not have believed it possible that the taste or the pride of Miss Fairfax could endure such society and friendship as the Vicarage had to offer.

»She is a riddle, quite a riddle!« said she.—»To chuse to remain here month after month, under privations of every sort! And now to chuse the mortification of Mrs Elton's notice and the penury of her conversation, rather than return to the superior companions who have always loved her with such real, generous affection.«

Jane had come to Highbury professedly for three months; the Campbells were gone to Ireland for three months; but now the Campbells had promised their daughter to stay at least till Midsummer, and fresh invitations had arrived for her to join them there. According to Miss Bates—it all came from her—Mrs Dixon had written most pressingly. Would Jane but go, means were to be found, servants sent, friends contrived—no travelling difficulty allowed to exist; but still she had declined it!

»She must have some motive, more powerful than appears, for refusing this invitation,« was Emma's conclusion. »She must be under some sort of penance, inflicted either by the Campbells or herself. There is great fear, great caution, great resolution somewhere.—She is not to be with the Dixons. The decree is issued by somebody. But why must she consent to be with the Eltons?—Here is quite a separate puzzle.«

Upon her speaking her wonder aloud on that part of the subject, before the few who knew her opinion of Mrs Elton, Mrs Weston ventured this apology for Jane.

»We cannot suppose that she has any great enjoyment at the Vicarage, my dear Emma—but it is better than being always at home. Her aunt is a good creature, but, as a constant companion, must be very tiresome. We must consider what Miss Fairfax quits, before we condemn her taste for what she goes to.«

»You are right, Mrs Weston,« said Mr Knightley warmly, »Miss Fairfax is as capable as any of us of forming a just opinion of Mrs Elton. Could she have chosen with whom to associate, she would not have chosen her. But (with a reproachful smile at Emma) she receives attentions from Mrs Elton, which nobody else pays her.«

Emma felt that Mrs Weston was giving her a momentary glance; and she was herself struck by his warmth. With a faint blush, she presently replied,

»Such attentions as Mrs Elton's, I should have imagined, would rather disgust than gratify Miss Fairfax. Mrs Elton's invitations I should have imagined any thing but inviting.«

»I should not wonder,« said Mrs Weston, »if Miss Fairfax were to have been drawn on beyond her own inclination, by her aunt's eagerness in accepting Mrs Elton's civilities for her. Poor Miss Bates may very likely have committed her niece and hurried her into a greater appearance of intimacy than her own good sense would have dictated, in spite of the very natural wish of a little change.«

Both felt rather anxious to hear him speak again; and after a few minutes silence, he said,

»Another thing must be taken into consideration too—Mrs Elton does not talk to Miss Fairfax as she speaks of her. We all know the difference between the pronouns he or she and thou, the plainest spoken amongst us; we all feel the influence of a something beyond common civility in our personal intercourse with each other—a something more early implanted. We cannot give any body the disagreeable hints that we may have been very full of the hour before. We feel things differently. And besides the operation of this, as a general principle, you may be sure that Miss Fairfax awes Mrs Elton by her superiority both of mind and manner; and that, face to face, Mrs Elton treats her with all the respect which she has a claim to. Such a woman as Jane Fairfax probably never fell in Mrs Elton's way before—and no degree of vanity can prevent her acknowledging her own comparative littleness in action, if not in consciousness.«

»I know how highly you think of Jane Fairfax,« said Emma. Little Henry was in her thoughts, and a mixture of alarm and delicacy made her irresolute what else to say.

»Yes,« he replied, »any body may know how highly I think of her.«

»And yet,« said Emma, beginning hastily and with an arch look, but soon stopping—it was better, however, to know the worst at once—she hurried on—»And yet, perhaps, you may hardly be aware yourself how highly it is. The extent of your admiration may take you by surprize some day or other.«

Mr Knightley was hard at work upon the lower buttons of his thick leather gaiters, and either the exertion of getting them together, or some other cause, brought the colour into his face, as he answered,

»Oh! are you there?—But you are miserably behindhand. Mr Cole gave me a hint of it six weeks ago.«

He stopped.—Emma felt her foot pressed by Mrs Weston, and did not herself know what to think. In a moment he went on—

»That will never be, however, I can assure you. Miss Fairfax, I dare say, would not have me if I were to ask her—and I am very sure I shall never ask her.«

Emma returned her friend's pressure with interest; and was pleased enough to exclaim,

»You are not vain, Mr Knightley. I will say that for you.«

He seemed hardly to hear her; he was thoughtful—and in a manner which shewed him not pleased, soon afterwards said,

»So you have been settling that I should marry Jane Fairfax?«

»No indeed I have not. You have scolded me too much for match-making, for me to presume to take such a liberty with you. What I said just now, meant nothing. One says those sort of things, of course, without any idea of a serious meaning. Oh! no, upon my word I have not the smallest wish for your marrying Jane Fairfax or Jane any body. You would not come in and sit with us in this comfortable way, if you were married.«

Mr Knightley was thoughtful again. The result of his reverie was, »No, Emma, I do not think the extent of my admiration for her will ever take me by surprize.—I never had a thought of her in that way, I assure you.« And soon afterwards, »Jane Fairfax is a very charming young woman—but not even Jane Fairfax is perfect. She has a fault. She has not the open temper which a man would wish for in a wife.«

Emma could not but rejoice to hear that she had a fault. »Well,« said she, »and you soon silenced Mr Cole, I suppose?«

»Yes, very soon. He gave me a quiet hint; I told him he was mistaken; he asked my pardon and said no more. Cole does not want to be wiser or wittier than his neighbours.«

»In that respect how unlike dear Mrs Elton, who wants to be wiser and wittier than all the world! I wonder how she speaks of the Coles—what she calls them! How can she find any appellation for them, deep enough in familiar vulgarity? She calls you, Knightley—what can she do for Mr Cole? And so I am not to be surprized that Jane Fairfax accepts her civilities and consents to be with her. Mrs Weston, your argument weighs most with me. I can much more readily enter into the temptation of getting away from Miss Bates, than I can believe in the triumph of Miss Fairfax's mind over Mrs Elton. I have no faith in Mrs Elton's acknowledging herself the inferior in thought, word, or deed; or in her being under any restraint beyond her own scanty rule of good-breeding. I cannot imagine that she will not be continually insulting her visitor with praise, encouragement, and offers of service; that she will not be continually detailing her magnificent intentions, from the procuring her a permanent situation to the including her in those delightful exploring parties which are to take place in the barouche-landau.«

»Jane Fairfax has feeling,« said Mr Knightley—»I do not accuse her of want of feeling. Her sensibilities, I suspect, are strong—and her temper excellent in its power of forbearance, patience, self-control; but it wants openness. She is reserved, more reserved, I think, than she used to be—And I love an open temper. No—till Cole alluded to my supposed attachment, it had never entered my head. I saw Jane Fairfax and conversed with her, with admiration and pleasure always—but with no thought beyond.«

»Well, Mrs Weston,« said Emma triumphantly when he left them, »what do you say now to Mr Knightley's marrying Jane Fairfax?«

»Why, really, dear Emma, I say that he is so very much occupied by the idea of not being in love with her, that I should not wonder if it were to end in his being so at last. Do not beat me.«